% !TEX root = Vicente_Calisto_Resume_EN.tex
%-------------------------------------------------------------------------------
% Layout and macros for the resume. Content lives in resume.tex.
%
% Expects \resumelang to be defined as `en` or `es` before this file is input
% (the entry points Vicente_Calisto_Resume_EN.tex / _ES.tex do this).
%
% Based on Jake Gutierrez's resume template, itself based on
% https://github.com/sb2nov/resume. Both are MIT-licensed; their notice:
%
%   Copyright (c) Jake Gutierrez, Sourabh Bajaj
%   Permission is hereby granted, free of charge, to any person obtaining a copy
%   of this software and associated documentation files (the "Software"), to deal
%   in the Software without restriction, including without limitation the rights
%   to use, copy, modify, merge, publish, distribute, sublicense, and/or sell
%   copies of the Software, and to permit persons to whom the Software is
%   furnished to do so, subject to the following conditions: The above copyright
%   notice and this permission notice shall be included in all copies or
%   substantial portions of the Software.
%   THE SOFTWARE IS PROVIDED "AS IS", WITHOUT WARRANTY OF ANY KIND, EXPRESS OR
%   IMPLIED, INCLUDING BUT NOT LIMITED TO THE WARRANTIES OF MERCHANTABILITY,
%   FITNESS FOR A PARTICULAR PURPOSE AND NONINFRINGEMENT. IN NO EVENT SHALL THE
%   AUTHORS OR COPYRIGHT HOLDERS BE LIABLE FOR ANY CLAIM, DAMAGES OR OTHER
%   LIABILITY, WHETHER IN AN ACTION OF CONTRACT, TORT OR OTHERWISE, ARISING FROM,
%   OUT OF OR IN CONNECTION WITH THE SOFTWARE OR THE USE OR OTHER DEALINGS IN THE
%   SOFTWARE.
%-------------------------------------------------------------------------------

\documentclass[letterpaper,11pt]{article}

%----------LANGUAGE----------
\providecommand{\resumelang}{en}
\newif\ifspanish
\makeatletter
\edef\@tempa{\resumelang}\def\@tempb{es}
\ifx\@tempa\@tempb \spanishtrue \fi
\makeatother

% \tr{English text}{Spanish text}: prints the text for the current language.
% Every piece of prose goes through it, so both languages are edited side by side.
\newcommand{\tr}[2]{\ifspanish#2\else#1\fi}

%----------ENGINE / FONTS----------
\usepackage{iftex}
\usepackage[T1]{fontenc}
\usepackage{lmodern}
\ifPDFTeX
  \usepackage[utf8]{inputenc}
  \input{glyphtounicode}
  \pdfgentounicode=1 % machine-readable text, for ATS parsers
\fi

\ifspanish
  \usepackage[english,spanish,es-noshorthands,es-nolayout,es-nodecimaldot]{babel}
\else
  \usepackage[spanish,english]{babel}
\fi

%----------PACKAGES----------
\usepackage[empty]{fullpage}
\usepackage{titlesec}
\usepackage{xcolor}
\usepackage{enumitem}
\usepackage{tabularx}
\usepackage[hidelinks]{hyperref}

% PDF metadata: shown in viewers' title bars and document properties, and read by ATS systems
\hypersetup{
  pdfauthor={Vicente Calisto},
  pdftitle={\tr{Vicente Calisto - Resume}{Vicente Calisto - Curr\'iculum}},
  pdfsubject={\tr{Senior Software Engineer, Database Engines (C++)}
                 {Ingeniero de Software Senior, Motores de Bases de Datos (C++)}},
  pdfkeywords={C++, C++14, C++17, C++20, database engines, database internals, query engine, query optimization,
    vectorized execution, columnar database, Apache Arrow, Parquet, SQL, SPARQL, graph database,
    vector search, HNSW, WebAssembly, Emscripten, Python, TypeScript, MillenniumDB, ColibriDB},
  pdflang={\tr{en-US}{es-CL}},
  pdfdisplaydoctitle=true, % viewers show the title instead of the file name
  pdfcreator={LaTeX},
}

%----------PAGE----------
\pagestyle{empty} % no header, footer or page number

\addtolength{\oddsidemargin}{-0.55in}
\addtolength{\evensidemargin}{-0.55in}
\addtolength{\textwidth}{1.1in}
\addtolength{\topmargin}{-.6in}
\addtolength{\textheight}{1.15in}

\urlstyle{same}
\raggedbottom
\raggedright
\setlength{\tabcolsep}{0in}

% Section title, its rule, and fixed space above and below (5pt each)
\titleformat{\section}{\scshape\raggedright\large}{}{0em}{}[\color{black}\titlerule]
\titlespacing*{\section}{0pt}{5pt}{5pt}

%----------SPACING----------
% All vertical spacing comes from these list settings; no \vspace hacks, so
% every gap of the same kind is the same size.
%   entries (jobs, projects): 0.1in indent, 2.5pt between entries
%   one-line entries (papers, skills): same indent, no extra space between
%   bullets under an entry: 1em indent, 1pt above the first one, none between
\newcommand{\resumeSubHeadingListStart}{%
  \begin{itemize}[leftmargin=0.1in, label={}, topsep=0pt, partopsep=0pt, parsep=0pt, itemsep=2.5pt]}
\newcommand{\resumeSubHeadingListEnd}{\end{itemize}}
\newcommand{\resumeCompactListStart}{%
  \begin{itemize}[leftmargin=0.1in, label={}, topsep=0pt, partopsep=0pt, parsep=0pt, itemsep=0pt]}
\newcommand{\resumeCompactListEnd}{\end{itemize}}
\newcommand{\resumeItemListStart}{%
  \begin{itemize}[leftmargin=1em, labelsep=0.4em, topsep=1pt, partopsep=0pt, parsep=0pt, itemsep=0pt,
                  label=\textbullet, before=\small]}
\newcommand{\resumeItemListEnd}{\end{itemize}}

%----------MACROS----------
\newcommand{\resumeItem}[1]{\item #1}

% {Organization}{Location}{Role}{Dates}
\newcommand{\resumeSubheading}[4]{%
  \item \textbf{#1}\hfill{\small #2}\par
  {\small\itshape #3\hfill #4\par}}

% {Second role at the same organization}{Dates}
\newcommand{\resumeSubSubheading}[2]{%
  \item {\small\itshape #1\hfill #2\par}}

% {Title | description}{Right-aligned text}
\newcommand{\resumeProjectHeading}[2]{%
  \item {\small #1\hfill #2\par}}

% {DOI}{Title}{Venue}: title, then a small underlined DOI link to the paper
\newcommand{\resumePublication}[3]{%
  \item {\small #2 {\footnotesize\href{https://doi.org/#1}{\underline{DOI}}}\hfill #3\par}}
